\section{\personFloer{}流线的 \texorpdfstring{$C^0$}{C⁰} 界}
\label{sec:c0-bounds-floer-trajectories}

\stfhIndexRangeStart{I001}{00-c0-jie}{$C^0$ 界}
若 $u$ 满足式~\eqref{eq:perturbed-cauchy-riemann-conormal}，则与一般梯度流方程的情形一样，$u$ 的能量可重写为
\begin{equation}
\begin{split}
  E_{(H,J)}(u)
  &=\int\!\!\int
    \left\lVert\frac{\partial u}{\partial t}-X_H(u)\right\rVert_{J_t}^2
    dt\,d\tau\\
  &=\mathcal A_H(u(-\infty))-\mathcal A_H(u(+\infty)).
\end{split}
\label{eq:floer-energy-action-drop}
\end{equation}
现在推导这一恒等式的若干结论。用 $\widetilde{\mathcal M}(H,J;z_-,z_+)$ 表示 $\widetilde{\mathcal M}(H,J)$ 中由满足\emph{固定}渐近条件
\[
  u(-\infty)=z_-,\qquad u(\infty)=z_+
\]
的 $u$ 所成的子集。

\begin{enumerate}[label=(\arabic*)]
  \item 若 $\widetilde{\mathcal M}(H,J;z_-,z_+)\ne\varnothing$，则
  \[
    \mathcal A_H(z_-)-\mathcal A_H(z_+)\ge0,
  \]
  且等号仅在 $z_-=z_+$ 并且 $u$ 静止时成立。
  \item 能量 $E_{(H,J)}(u)$ 只依赖于渐近哈密顿轨道 $u(\pm\infty)$ 的作用量值。对所有
  \[
    u\in\widetilde{\mathcal M}(H,J;z_-,z_+),
  \]
  我们有\emph{\termAprioriEnergyBound{}}
  \begin{equation}
    E_{(H,J)}(u)
    =\mathcal A_H(z_-)-\mathcal A_H(z_+)<\infty.
    \label{eq:apriori-energy-bound-asymptotes}
  \end{equation}
  \item 当 $z_-\ne z_+$ 时，$\tau$-平移在 $\widetilde{\mathcal M}(H,J;z_-,z_+)$ 上给出一个自然的自由 $\mathbb R$-作用。我们记
  \[
    \mathcal M(H,J;z_-,z_+)
    =\widetilde{\mathcal M}(H,J;z_-,z_+)/\mathbb R.
  \]
\end{enumerate}

在第一卷第 2 部分中，我们已经在真正的 $J$-全纯方程背景下，发展了研究模空间所需的分析。例如，\termBubblingOffAnalysis{}同样适用于这些受扰 $J$-全纯方程。但是，对 $T^*N$ 中像 $\nu^*S_0,\nu^*S_1$ 这样的正合拉格朗日对子，非常值球面或圆盘的冒泡不可能发生。\emph{只要}交集是紧的，并且对解 $u$ 有某些 $C^0$-估计，那么在\emph{紧}\termLagrangianSubmanifold{}情形下发展的一切必要\termFredholmProperty{}与紧性性质，现在都可用于当前的非紧目标空间。在非紧\termSymplecticManifold{}中的\emph{非紧}\termLagrangianSubmanifold{}上建立\personFloer{}理论时，这个 $C^0$-估计是关键一步。为了建立哈密顿量或子流形 $S$ 在形变下的某种不变性，我们还需考虑式~\eqref{eq:perturbed-cauchy-riemann-conormal} 的参数化版本。

现在把讨论分成两种情形：一种用于构造\termFloerBoundaryMap{}，另一种用于构造\termChainMap{}，特别是在子流形同痕 $\{S_s\}_{0\le s\le1}$ 导致边界 $\nu^*S_s$ 改变时。

在第一种情形中，我们研究只涉及固定边界的方程~\eqref{eq:perturbed-cauchy-riemann-conormal}。由于所考虑的哈密顿量 $H$ 使 $X_H$ 紧支，在这种情形下可以使用\termStrongMaximumPrinciple{}。下面对此作出说明。

先回顾\personGromov{} (Gr85) 意义下 $J$-凸性的概念。

\begin{definition}
\label{def:contact-type-boundary}
设 $(Z,\omega)$ 是带边界 $Q=\partial Z$ 的\termSymplecticManifold{}。若存在 $Q$ 上的接触形式 $\lambda$\footnote{译者注：原书此处写为“contact from $\lambda$ on $Z$”；`from` 应为 `form`，且接触形式应定义在 $Q$ 上。原书第 (2) 项又将边界写成未定义的 $\partial W$；译文改为 $\partial Z$。}，使得
\begin{enumerate}[label=(\arabic*)]
  \item $\omega|_Q=d\lambda$ 且 $\xi=\ker\lambda$；
  \item 由 $\lambda\wedge(d\lambda)^{n-1}$ 定义的 $Q$ 的定向，与 $Q=\partial Z$ 的边界定向一致。
\end{enumerate}
则称 $(Z,\omega)$ 具有\termContactTypeBoundary{}。
\end{definition}

\begin{definition}
\label{def:j-convex-boundary}
设 $(Z,\omega,J)$ 是一个具有接触型边界的殆\personKahler{}结构。若 $\partial Z$ 上存在接触形式 $\lambda$，使得二形式 $d\lambda(\mathord\cdot,J\mathord\cdot)$ 在 $\xi$ 上正定，则称 $\partial Z$ 是 $J$-凸的。
\end{definition}

有了这些定义，便得到如下结论。

\begin{proposition}
\label{prop:j-convex-maximum-principle}
设 $(Z,\omega,J)$ 是一个殆\personKahler{}结构，其边界 $\partial Z$ 是 $J$-凸的。则 $Z$ 中的任意 $J$-全纯曲线 $u$ 都不能在 $u$ 的定义域内点从内侧接触 $\partial Z$。更精确地说，不存在点 $z_0\in\operatorname{Int}\operatorname{Dom}u$ 及其开邻域 $U\subset\operatorname{Dom}u$，使得
\[
  u(z_0)\in\partial Z,
  \qquad
  u(U\setminus\{z_0\})\subset\operatorname{Int}Z.
\]
\end{proposition}

\begin{proof}
边界的 $J$-凸性意味着距离函数
\[
  \psi(x):=\operatorname{dist}(x,\partial Z)^2
\]
成为一个\termPlurisubharmonicFunction{}，也就是说，它满足 $d(d\psi\circ J)$ 正定这一判别条件。

反设存在这样的内点 $z_0$，并设 $U$ 是相应的开邻域。那么函数 $\psi\circ u$ 成为次调和函数。这与最大值原理矛盾，从而完成证明。\footnote{译者注：按原书所写 $\psi=\operatorname{dist}(\mathord\cdot,\partial Z)^2\ge0$，内部接触点给出的是局部最小值 0；原书却称 $\psi\circ u$ 次调和并调用最大值原理。这里存在符号或最值原理方向的缺口，译文不自行重构证明。}
\end{proof}

下面给出主要的能量界和 $C^0$ 界。在 $T^*N$ 中写
\[
  u(\tau,t)=(q(\tau,t),p(\tau,t)),
\]
并用 $|\mathord\cdot|_g$ 表示由 $N$ 上度量 $g$ 诱导的 $T_{q(\tau,t)}^*N$ 上的范数。

\begin{theorem}
\label{thm:cotangent-floer-c0-bound}
假设 $S_0,S_1\subset N$ 横截相交。则存在常数 $C=C(H,J,S_0,S_1)>0$，使得式~\eqref{eq:perturbed-cauchy-riemann-conormal} 的每个像有界的有限能量解 $u$ 都满足\footnote{译者注：原书引入了常数 $C$，但下式及后文只使用 $R_0$；下式中取上确界的范围 $\Theta$ 也未定义，从 $u$ 的定义域看应指 $\mathbb R\times[0,1]$。译文照录这两个原书记号。}
\begin{equation}
  \sup_{(\tau,t)\in\Theta}|p(\tau,t)|_g\le R_0,
  \label{eq:cotangent-fiber-c0-bound}
\end{equation}
其中 $R_0>0$ 由下式给出：
\[
  R_0
  :=\inf\{R\in\mathbb R_+
    \mid T^*N\setminus D_R
    \subset T^*N\setminus(\operatorname{Supp}H\cup\operatorname{Supp}J)\}.
\]
\end{theorem}

\begin{proof}
由有限能量假设，我们有\footnote{译者注：原书下式的\personHamilton{}项写作 $X_K(u)$，但此处的哈密顿量为 $H$；译文改为 $X_H(u)$。}
\[
  \int_{-\infty}^{\infty}\!\int_0^1
  \left(
    \left\lVert\frac{\partial u}{\partial\tau}\right\rVert_{J_t}^2
    +\left\lVert\frac{\partial u}{\partial t}-X_H(u)\right\rVert_{J_t}^2
  \right)dt\,d\tau<\infty.
\]
再由\cref{lem:finite-energy-subsequence-convergence}以及 $u$ 的\termExponentialDecayEstimate{}（在稍有不同的背景下，证明见第 14.1 节；指数估计的一般框架见附录 B.1），得到一致收敛
\[
  \lim_{\tau\to-\infty}u(\tau)=z_-,
  \qquad
  \lim_{\tau\to+\infty}u(\tau)=z_+.
\]
这里 $z_-$ 和 $z_+$ 分别满足\footnote{译者注：原书把 $z_+$ 的极限也写成 $\tau\to-\infty$，且式~(12.3.15) 的右边写作 $X_H(t,z)$。译文分别改为 $\tau\to+\infty$ 与 $X_H(t,z_\pm)$。}
\begin{equation}
  \dot z_\pm=X_H(t,z_\pm),
  \qquad
  z_\pm(0)\in\nu^*S_0,
  \quad z_\pm(1)\in\nu^*S_1.
  \label{eq:asymptotic-hamiltonian-paths}
\end{equation}
于是必有下列两种情形之一：写 $u(\tau,t)=(q(\tau,t),p(\tau,t))$，则 $\sup_{(\tau,t)}|p(\tau,t)|_g$ 要么在 $\tau=\pm\infty$ 处取到，要么在某点 $(\tau_0,t_0)\in\mathbb R\times[0,1]$ 处取到。

由 $X_H$ 紧支，从而 $\phi_H^1(\nu^*S_0)\cap\nu^*S_1$ 是紧的，不难推导出哈密顿路径 $z^\alpha$ 的 $C^0$-估计。\footnote{译者注：原书此处两次使用 $z^\alpha$，但 $\alpha$ 未定义；前文定义的渐近道路是 $z_-$ 与 $z_+$。译文照录原记号。}因此，一旦得到路径 $z^\alpha$ 的 $C^0$-估计，第一种情形的证明就完成了。所以只需考虑第二种情形。证明下述命题就足够了。

\begin{proposition}
\label{prop:maximum-principle-disc-bundle}
假设对某个 $R>0$ 有
\[
  \bigl(\phi_H^1(\nu^*S_0)\cap\nu^*S_1\bigr)
  \cup\operatorname{supp}H\cup\operatorname{supp}J
  \subset D_R,
\]
其中 $D_R\subset T^*N$ 是\termDiscBundle{}
\[
  D_R:=\{p\in T^*N\mid |p|_g\le R\}.
\]
则对式~\eqref{eq:perturbed-cauchy-riemann-conormal} 的所有解 $u$，都有 $u(\tau_0,t_0)\in D_R$。
\end{proposition}

\begin{proof}
反设 $u(\tau_0,t_0)\notin D_R$。不妨设 $u(\tau_0,t_0)\in\partial D_{R'}$，其中 $R'>R$，并且对某个 $\epsilon>0$，$(\tau_0,t_0)$ 在 $\mathbb R\times[0,1]$ 中有开邻域 $B_\epsilon$，使得
\[
  u(B_\epsilon)
  \subset T^*N\setminus D_R
  \subset T^*N\setminus(\operatorname{Supp}H\cup\operatorname{Supp}J).
\]
因此 $u$ 在 $B_\epsilon$ 上满足
\[
  \frac{\partial u}{\partial\tau}
  +J_g\frac{\partial u}{\partial t}=0.
\]
分别考虑两种情形：$(\tau_0,t_0)\in\operatorname{Int}(\mathbb R\times[0,1])$ 和 $(\tau_0,t_0)\in\partial(\mathbb R\times[0,1])$。

\cref{prop:j-convex-maximum-principle}立即排除了 $(\tau_0,t_0)\in\operatorname{Int}(\mathbb R\times[0,1])$ 的可能性。

现在考虑 $(\tau_0,t_0)\in\partial(\mathbb R\times[0,1])$ 的情形，即 $t_0=0,1$。$t_0=0$ 的情形与之相同，故只考虑 $t_0=1$。考虑边界曲线
\[
  \tau\longmapsto u(\tau,1)=(q(\tau,1),p(\tau,1)),
\]
它在 $(\tau_0,1)$ 处与 $\partial D_{R'}\cap\nu^*S$ 相切。由于 $\nu^*S\cap\partial D_{R'}$ 是 $\partial D_{R'}$ 中的\termLegendrianSubmanifold{}，该曲线在 $(\tau_0,1)$ 处与 $\partial D_{R'}$ 的\termContactDistribution{}相切。\footnote{译者注：原书在本句把前后一贯使用的 $\partial D_{R'}$ 写成 $\partial D_r$，但紧随的陈列式又恢复为 $\partial D_{R'}$；译文统一为后者。}这个接触分布为
\[
  \left\{\xi\in T(\partial D_{R'})
    \mathrel{}\middle|\mathrel{}\xi\perp J_g\frac{\partial}{\partial r}\right\},
\]
其中 $\frac{\partial}{\partial r}$ 是 $T^*N$ 上的\termRadialVectorField{}。

由于 $u$ 是 $J_g$-全纯的，而接触分布 $\xi$ 是 $J_g$-不变的，$u$ 在 $(\tau_0,1)$ 处也与接触分布相切。特别地，
\[
  \left.\frac{\partial}{\partial t}|p(\tau,t)|^2\right|_{(\tau_0,1)}=0.
\]
然而，对 $B_\epsilon(\tau_0,1)$ 上（相对于 $J_g$）的次调和函数
\[
  (\tau,t)\longmapsto |p(\tau,t)|^2
\]
使用\termStrongMaximumPrinciple{}，就会得到矛盾。这是因为，如有必要，可以把 $\epsilon$ 选得稍大，从而假设
\[
  \left.\operatorname{Image}u\right|_{B_\epsilon(\tau_0,1)}
  \not\subset\partial D_{R'}.
\]
这就完成了该命题的断言，从而也完成了\cref{thm:cotangent-floer-c0-bound}的证明。
\end{proof}

\begin{exercise}
\label{ex:fiber-norm-subharmonic}
证明函数 $(\tau,t)\longmapsto|p(\tau,t)|^2$ 是次调和函数。
\end{exercise}

这就完成了定理的证明。
\end{proof}
\stfhIndexRangeEnd{I001}{00-c0-jie}{$C^0$ 界}
